% Prior work (deck, slide "Prior Work"), second style (the look of the ScaleQsim paper).
% What the figure shows: one table of the existing implementations of Pauli propagation and of SectorQ, one row
% per implementation. Under each name stands where it is from and its year: (Source, Year).
% The rows are grouped by where the implementation runs (first column, a cell that spans its group, with a
% small pictogram): one GPU -> one CPU node -> many CPU nodes -> many GPUs. For every implementation the three
% right columns give what this work adds: when terms move between processes, whether a run can hold more terms
% than the GPU memory, and how the terms are split over the processes.
% Message (one sentence): existing implementations stay on one device, run on CPUs, or move terms between
% all GPUs at every gate; SectorQ runs on many GPUs, moves terms only at window boundaries, holds more terms
% than GPU memory, and chooses the split at run time.
% Sources of the entries: README 1.3, 2.4, 2.5, 3.2 to 3.4, 4.2. A dash is a property that the sources used
% here do not state. No measured numbers.
% (Source, Year), from the citation files of the code that was run (pauliprop-sectorq-impl/third_party):
%   SPD                  github.com/Quantinuum/spd, no paper named; version used: commit of 2026-09-28
%   PauliPropagation.jl  Rudolph et al., PRX Quantum 7, 032001 (2026)
%   ppvm                 github.com/QuEraComputing/ppvm, no paper named; version used: commit of 2026-10-01
%   Qiskit pauli-prop    software citation of the Qiskit addons team, 2025
%   ORQA                 Broers et al. (RIKEN), arXiv 2506.13241 (2025)
%   MonoProp             Algorithmiq; preferred citation Miller et al., arXiv 2503.18939 (2025)
%   cuPauliProp          NVIDIA cuQuantum, no paper; version used: v0.5.0 of cuQuantum 26.09 (2026)
% Colors: blue = header row, gray = the group of a row, peach/teal/yellow/lavender boxes = nodes, khaki boxes =
% GPUs, yellow = this work, red = the key value of the figure.
% Sizes: the canvas is 17.6 x 8.86 cm and is shown 1060 px wide on the slide; names and cells are 10 pt, the
% (Source, Year) line 8.6 pt italic, the row of this work 10.5 pt and its key value 11.5 pt.
\begin{tikzpicture}[x=1cm, y=1cm, font=\rmfamily\bfseries\fontsize{10}{11.2}\selectfont]
  \tikzset{
    every node/.append style={execute at begin node=\frenchspacing},     % no wider space after a colon
    % a cell: an (empty or centered) box of a given width; the text of the implementation column is put on top
    c/.style={draw=black, fill=white, line width=0.5pt, anchor=north west, inner sep=0pt, align=center,
              font=\rmfamily\bfseries\fontsize{10}{11.2}\selectfont, minimum height=0.76cm},
    hd/.style={c, fill=thblue, text=white, minimum height=0.88cm},
    bd/.style={c, minimum height=0.50cm, font=\rmfamily\bfseries\fontsize{10.5}{12}\selectfont},
    grp/.style={c, fill=pan2},
    hl/.style={c, fill=actc!55, minimum height=0.96cm, font=\rmfamily\bfseries\fontsize{10.5}{12}\selectfont},
    nm/.style={anchor=west, inner sep=0pt, align=left, font=\rmfamily\bfseries\fontsize{10}{10.6}\selectfont},
    src/.style={font=\rmfamily\bfseries\itshape\fontsize{8.6}{9.4}\selectfont, text=black!78},
    gname/.style={font=\rmfamily\bfseries\fontsize{10.5}{11.6}\selectfont, inner sep=0pt, align=center},
    kv/.style={font=\rmfamily\bfseries\itshape\fontsize{11.5}{12.4}\selectfont, text=dred, inner sep=0pt, align=center},
    pic/.style={draw=black!85, fill=white, line width=0.6pt, inner sep=0pt, minimum width=0.36cm, minimum height=0.36cm},
  }
  \def\ns{\textcolor{black!55}{---}}
  % columns: left edge and width (the table is 17.40 cm wide, from 0.10 to 17.50)
  \def\cxa{0.10}  \def\cwa{2.40}     % where it runs (the group)
  \def\cxb{2.50}  \def\cwb{4.95}     % implementation (Source, Year)
  \def\cxc{7.45}  \def\cwc{4.20}     % when terms move between processes
  \def\cxd{11.65} \def\cwd{3.25}     % holds more terms than GPU memory
  \def\cxe{14.90} \def\cwe{2.60}     % how the terms are split
  \def\cxr{17.50}                    % right edge of the table
  \def\cgx{1.30}                     % center of the group column
  % rows (top edge): band -0.76, header -1.26, implementations -2.14 ... -6.70 (0.76 cm each), this work -7.46
  % (0.96 cm), end of the table -8.42
  % =====================================================================================================
  % panel and title
  % =====================================================================================================
  \draw[pan] (0,0) rectangle (17.60,-8.86);
  \node[ptitle, font=\rmfamily\bfseries\LARGE, anchor=base west] at (0.10,-0.53)
    {Existing Implementations of Pauli Propagation and SectorQ};
  % =====================================================================================================
  % band and header row
  % =====================================================================================================
  \node[bd, minimum width=7.35cm] at (\cxa,-0.76) {\strut Shared With Prior Work};
  \node[bd, fill=actc, minimum width=10.05cm] at (\cxc,-0.76) {\strut Added by This Work};
  \node[hd, minimum width=\cwa cm] at (\cxa,-1.26) {\strut Where It\\\strut Runs};
  \node[hd, minimum width=\cwb cm] at (\cxb,-1.26) {\strut Implementation\\\strut {\itshape (Source, Year)}};
  \node[hd, minimum width=\cwc cm] at (\cxc,-1.26) {\strut When Terms Move\\\strut Between Processes};
  \node[hd, minimum width=\cwd cm] at (\cxd,-1.26) {\strut Holds More Terms\\\strut Than GPU Memory};
  \node[hd, minimum width=\cwe cm] at (\cxe,-1.26) {\strut How the Terms\\\strut Are Split};
  % =====================================================================================================
  % the seven existing implementations: name, (source, year), and the three properties
  % =====================================================================================================
  \foreach \y/\n/\w/\a/\b/\s in {
      -2.14/{SPD}/{(Quantinuum, 2026)}/{\strut Never: One Device}/{\strut No}/{\strut\ns},
      -2.90/{PauliPropagation.jl}/{(PRX Quantum, 2026)}/{\strut Never: One Device}/{\strut\ns}/{\strut\ns},
      -3.66/{ppvm}/{(QuEra, 2026)}/{\strut Never: One Device}/{\strut\ns}/{\strut\ns},
      -4.42/{Qiskit pauli-prop}/{(Qiskit, 2025)}/{\strut Never: One Device}/{\strut\ns}/{\strut\ns},
      -5.18/{ORQA}/{(RIKEN, arXiv 2025)}/{\strut At Every Branching Gate}/{\strut\ns}/{\strut Fixed Function},
      -5.94/{MonoProp}/{(Algorithmiq, arXiv 2025)}/{\strut\ns}/{\strut\ns}/{\strut\ns},
      -6.70/{cuPauliProp Distributed Mode}/{(NVIDIA, 2026)}/{\strut At Every Gate}/{\strut No}/{\strut Not Public}}
  {
    \node[c, minimum width=\cwb cm] at (\cxb,\y) {};
    \node[nm] at ({\cxb+0.16},{\y-0.25}) {\strut\n};
    \node[nm, src] at ({\cxb+0.16},{\y-0.565}) {\strut\w};
    \node[c, minimum width=\cwc cm] at (\cxc,\y) {\a};
    \node[c, minimum width=\cwd cm] at (\cxd,\y) {\b};
    \node[c, minimum width=\cwe cm] at (\cxe,\y) {\s};
  }
  % =====================================================================================================
  % this work
  % =====================================================================================================
  \node[hl, minimum width=\cwb cm] at (\cxb,-7.46) {};
  \node[nm, font=\rmfamily\bfseries\fontsize{11}{12}\selectfont] at ({\cxb+0.16},-7.80) {\strut SectorQ};
  \node[nm, src, text=black] at ({\cxb+0.16},-8.15) {\strut (This Work)};
  \node[hl, minimum width=\cwc cm] at (\cxc,-7.46) {};
  \node[hl, minimum width=\cwd cm] at (\cxd,-7.46) {\strut Yes: Host Memory};
  \node[hl, minimum width=\cwe cm] at (\cxe,-7.46) {};
  \node[kv] at ({\cxc+0.5*\cwc},-7.955) {\strut Only at Window\\\strut Boundaries};
  \node[gname] at ({\cxe+0.5*\cwe},-7.955) {\strut Chosen at\\\strut Run Time};
  % =====================================================================================================
  % the four groups: a cell that spans the rows of the group, a pictogram and the name
  % =====================================================================================================
  \node[grp, minimum width=\cwa cm, minimum height=0.76cm] at (\cxa,-2.14) {};
  \node[grp, minimum width=\cwa cm, minimum height=2.28cm] at (\cxa,-2.90) {};
  \node[grp, minimum width=\cwa cm, minimum height=1.52cm] at (\cxa,-5.18) {};
  \node[grp, minimum width=\cwa cm, minimum height=1.72cm] at (\cxa,-6.70) {};
  % one GPU: one GPU box at the left of the name
  \node[pic, gpu] at (0.47,-2.52) {};
  \node[gname] at (1.58,-2.52) {\strut One GPU};
  % one CPU node: one node box
  \node[pic, node0] at (\cgx,-3.72) {};
  \node[gname] at (\cgx,-4.36) {\strut One CPU\\\strut Node};
  % many CPU nodes: four node boxes
  \foreach \i/\st in {0/node0, 1/node1, 2/node2, 3/node3}
    \node[pic, \st] at ({\cgx+0.46*(\i-1.5)},-5.50) {};
  \node[gname] at (\cgx,-6.18) {\strut Many CPU\\\strut Nodes};
  % many GPUs: four GPU boxes
  \foreach \i in {0,1,2,3}
    \node[pic, gpu] at ({\cgx+0.46*(\i-1.5)},-7.26) {};
  \node[gname] at (\cgx,-7.86) {\strut Many GPUs};
  % heavier lines between the groups; the frame of this work
  \draw[line width=1.1pt] (\cxa,-2.90) -- (\cxr,-2.90);
  \draw[line width=1.1pt] (\cxa,-5.18) -- (\cxr,-5.18);
  \draw[line width=1.1pt] (\cxa,-6.70) -- (\cxr,-6.70);
  \draw[navy, line width=1.8pt] (\cxb,-7.46) rectangle (\cxr,-8.42);
  % footnotes
  \node[ann, font=\rmfamily\bfseries\itshape\fontsize{8.6}{9.4}\selectfont, anchor=west] at (0.12,-8.645)
    {\ns\ : Not Stated in the Sources Used Here.\quad Year: of the Paper or the Software Citation; Without One, of the Code Version Used.};
\end{tikzpicture}
